\section{Introduction}
One of quantum computing's most significant results is the development of Grover's
algorithm~\cite{bib-grover1996fast}. This algorithm performs an unstructured search
optimally~\cite{bib-zaika-grov-opt}. Various applications have
been proposed for such an algorithm, including the acceleration of solving NP-hard
problems~\cite{bib-williams-grover-np}.

Grover's algorithm uses a phase oracle, a quantum circuit that applies a phase of
$-1$ when the Boolean function it implements is true, to increase the probability
of returning a fulfilling assignment.

{\it Quantum Boolean circuits} are a similar class of circuits that map bitstrings of
qubits to a different mapping of bitstrings. Several method have been developed
to generate them and optimize them~\cite{bib-meuli-mult,bib-amy-rm,bib-meuli-esop}.
In particular, many of these methods optimize quantum Boolean circuits for T-count, the
total number of $T$ and $T^{\dagger}$ gates in the circuit, and for the T-depth,
the longest chain of of $T$ and $T^{\dagger}$ gates in the circuit.

Some of the most promising methods use {\it Boolean logic networks}, which are
representations of Boolean functions that depict its decompositions as a
directed acyclic graph where each of the nodes is a subfunction in the
decomposition~\cite{bib-meuli-xag,bib-lhrs-soeken}. Properties of these Boolean
logic networks have been demonstrated to correspond directly to values
of interest to quantum circuit optimization. Of note, are XOR AND Inverter
graphs (XAG), which are Boolean logic networks consisting of the XOR and AND
operations, with any combinations of their inputs and outputs negated. The
 {\it multiplicative complexity}, or the number of AND gates of such a circuit
is directly related to the T-count~\cite{bib-meuli-mult}, and minimizing the
number of such AND nodes is connected to a reduction in T-count. Additionally,
an XAG's {\it multiplicative depth}, or the largest number of AND nodes on a
path between the primary input and a primary output, can be directly related
to the T-count of an oracle~\cite{bib-meuli-ros}.

There are existing methods to optimize the T-count and T-depth in quantum Boolean
circuits by additionally optimizing the {\it multiplicative complexity}~\cite{bib-meuli-mult}
and the T-depth by optimizing the {\it multiplicative depth}~\cite{bib-meuli-ros}.
These take advantage of the fact that Boolean logic networks can be rewritten
to share redundant logic between multiple different mappings among its output
qubits

These can be used to create a {\it quantum oracle}, which is a quantum circuit
that applies the X operation to ancillas when the function it is acting as an
oracle for is true. This can then be modeled as a quantum Boolean circuit that
maps all input qubits to itself, except for the ancilla it acts on. Such a construction
can also be called a {\it single target gate}. We demonstrate that such a quantum oracle
can then be converted to a phase oracle in a straightforward manner. 

We found, however, that applying existing methods in such a straightforward manner
unnecessarily increases T-count and T-depth. In particular, implementing the oracle
as a single target gate uses more T-gates than would otherwise be needed. It is
generally more efficient to directly generate such phase oracles using phase gates
as much as possible. 

We demonstrate that any phase gate controlled by three or fewer variables can be implemented using
only phase gates such as the $T$ and $T^{\dagger}$ gates and CNOT gates (CNOT+T). This
gate set is notable because there exist methods to freely rearrange it for lower
T-depth. Additionally, for a multi-level logic network such as that generated by
abc~\cite{bib-mishchenko-lut}, the last layer can be implemented as Controlled-Z gates,
thus saving on T-depth for product terms of over three variables, and allowing for more
simplification by sum of paths calculations.

In this paper, we introduce adaptations of existing methods used to generate quantum
Boolean circuits in order to realize savings in T-count and T-depth in phase oracles.
We combine three or fewer qubit phase oracles, ESOP optimization and phase oracle synthesis,
and LHRS to further reduce T-count and T-depth in phase oracle generation.

The contribution of this paper is as follows:

\begin{itemize}
\item A method to generate controlled-Z gates of three variables or fewer as
  CNOT+T circuits using the Boolean Fourier transform
\item A method to generate the phase oracle as a Boolean logic network from the
  simplified ESOP expression, and then implementing its nodes as 2-input Boolean functions,
  with the last layer of nodes implemented by Controlled-Z gates from the library generated
  in the first step.
\end{itemize}

We compare our method to Qiskit's \texttt{PhaseOracle}~\cite{bib-phaseoracle} and the Caterpillar
library~\cite{bib-meuli-ros} and find that on average, our proposal reduces the T-count by 52\% of
each method, and the T-depth to 63\% of the original.

The rest of this paper is structured as follows: First Sec~\ref{Pre}
introduces some preliminary knowledge. In addition to the basics of
qubits and quantum circuits, this includes the basics of the
mathematics of phase polynomials and their mapping to phase oracles
and CNOT+T circuits. Sec.~\ref{Pro} details the proposal itself,
including how to generate a library of up to 3-qubit phase oracles,
and how to use it to generate phase oracles of 4 or more qubits.
In Sec.~\ref{Exp}, we test our results against existing methods
and analyze our findings. Finally, Sec.~\ref{Conc} concludes with
our findings and proposes some future research.

